\section{Related Work}
\label{submission:related_work}

\textbf{Weight-Space Model Merging.} Combining multiple neural networks in weight space has emerged as a promising approach for scaling model capabilities without additional training overhead. Early works such as Model Soups \cite{model_soups} demonstrated that simple averaging of fine-tuned weights from different training runs can enhance generalization and robustness. For multi-task setups, Task Arithmetic \cite{task_arithmetic} showed that taking parameter differences (task vectors) from fine-tuned task experts and adding them to a shared pre-trained base model effectively merges diverse task capabilities. However, parameter interference and conflicting parameter updates can severely degrade performance. To address this, TIES-Merging \cite{ties_merging} introduced a heuristic pipeline to resolve sign conflicts and prune redundant parameters. Other approaches, such as Git Re-Basin \cite{git_rebasin} and ZipIt! \cite{zipit}, utilize weight permutation and activation alignment to merge models from different initializations. While effective, these methods introduce significant geometric or combinatorial complexity.

\textbf{Adaptive and Test-Time Merging.} To avoid static merging scales, several works have proposed learning merging coefficients dynamically. Built upon the principles of unsupervised test-time adaptation via prediction entropy minimization \cite{tent}, AdaMerging \cite{AdaMerging_ICLR_2024} dynamically optimizes merging coefficients using a small calibration batch from the test-time distribution. In their work, Yang et al.\ \cite{AdaMerging_ICLR_2024} evaluate two specific configurations: (1) \emph{Task-wise AdaMerging}, which optimizes exactly $T$ coefficients (one global scalar per task), and (2) \emph{Layer-wise AdaMerging}, which optimizes $L \times T$ layer-specific scales. Representation Surgery \cite{representation_surgery} further extends this by aligning representational statistics across task experts during merging. Although adaptive merging achieves high in-distribution accuracies, it requires optimizing multiple scaling variables on tiny, unlabeled datasets, which exposes the system to transductive overfitting risks.

\textbf{Deconstructive and Regularization Analyses in ML.} Our work is deeply inspired by a growing body of machine learning research that systematically deconstructs over-engineered pipelines to isolate their true causal drivers. For instance, recent critiques of Sharpness-Aware Isotropic Merging (SAIM) \cite{saim_deconstruction} have shown that its complex coordinate-restricted optimization and SVD-based projections are largely redundant, and that standard global SAM fine-tuning is the primary contributor to performance. Similarly, analyses of FoldMerge \cite{convoluted_origami} have argued that learning continuous diffeomorphisms via normalizing flows introduces massive, unnecessary parameter overhead that can be outperformed by simpler, parameter-free alignment strategies. 

While prior work on the overfitting-optimizer paradox \cite{adamerging_paradox} indicated that layer-wise coefficients overfit to the test sample, the community has continued to use high-dimensional adaptation or seek increasingly complex objectives. We build upon and formalize these ideas by designing rigorous, empirical diagnostic controls---such as layer shuffling and spatial averaging---and presenting the first comparative deconstruction of the optimization dynamics of task-wise vs.\ layer-wise AdaMerging. This allows us to expose and mathematically explain the Spatial Averaging Paradox.
